%%%%%%%%%%%%%%%%%%%%%%%%%%%%%%%%%%%%%%%%%%%%%%%%%%%%%%%%%%%%%%%%%%%%%%%%%%%%%%%
%                       CARGA DE LA CLASE DE DOCUMENTO                        %
%                                                                             %
% Las opciones admisibles son:                                                %
%      12pt / 11pt            (tamaño del cuerpo de letra; no usar 10pt)      %
%                                                                             %
% catalan/spanish/english     (idioma principal del trabajo)                  %
%                                                                             % 
% french/italian/german...    (si necesitáis usar otro idioma adicional)      %
%                                                                             %
% listoffigures               (El documento incluye un Índice de figuras)     %
% listoftables                (El documento incluye un Índice de tablas)      %
% listofquadres               (El documento incluye un Índice de cuadros)     %
% listofalgorithms            (El documento incluye un Índice de algoritmos)  %
%                                                                             %
%%%%%%%%%%%%%%%%%%%%%%%%%%%%%%%%%%%%%%%%%%%%%%%%%%%%%%%%%%%%%%%%%%%%%%%%%%%%%%%

\documentclass[11pt,spanish,listoffigures,listoftables]{tfgetsinf}

%%%%%%%%%%%%%%%%%%%%%%%%%%%%%%%%%%%%%%%%%%%%%%%%%%%%%%%%%%%%%%%%%%%%%%%%%%%%%%%
%                     CODIFICACIÓN DEL ARCHIVO FUENTE                         %
%                                                                             %
%    Windows suele usar 'ansinew'                                             %
%    en Linux es posible que sea 'latin1' o 'latin9'                          %
%    Pero lo más recomendable es usar utf8 (unicode 8)                        %
%                                          (si vuestro editor lo permite)     % 
%%%%%%%%%%%%%%%%%%%%%%%%%%%%%%%%%%%%%%%%%%%%%%%%%%%%%%%%%%%%%%%%%%%%%%%%%%%%%%%

\usepackage[utf8]{inputenc} 


%%%%%%%%%%%%%%%%%%%%%%%%%%%%%%%%%%%%%%%%%%%%%%%%%%%%%%%%%%%%%%%%%%%%%%%%%%%%%%%
%                       OTROS PAQUETES Y DEFINICIONES                         %
%                                                                             %
%%%%%%%%%%%%%%%%%%%%%%%%%%%%%%%%%%%%%%%%%%%%%%%%%%%%%%%%%%%%%%%%%%%%%%%%%%%%%%%

\usepackage{glossaries}
\usepackage{textcomp}
\usepackage{booktabs}
\usepackage{float}
\usepackage{enumitem}
\usepackage{graphicx}
\usepackage{subcaption}
\usepackage{tabularx}
\usepackage{pgfplots}
\pgfplotsset{compat=1.18}

% Listings para código
\usepackage{listings}
\lstset{
  basicstyle=\ttfamily\small,
  frame=single,
  breaklines=true,
  postbreak=\mbox{\textcolor{gray}{$\hookrightarrow$}\space},
  columns=fullflexible,
  keepspaces=true,
  numbers=none
}
 
% Bibliografía
\usepackage[backend=biber,style=numeric,sorting=none]{biblatex}
\addbibresource{bibliografia.bib}

% Hyperref siempre al final
\usepackage{hyperref}
\hypersetup{
  colorlinks=true,
  linkcolor=black,
  urlcolor=cyan,
  citecolor=black
}
%%%%%%%%%%%%%%%%%%%%%%%%%%%%%%%%%%%%%%%%%%%%%%%%%%%%%%%%%%%%%%%%%%%%%%%%%%%%%%%
%                          DATOS DEL TRABAJO                                  %
%                                                                             %
% título alumno, titor y curso académico                                      %
%%%%%%%%%%%%%%%%%%%%%%%%%%%%%%%%%%%%%%%%%%%%%%%%%%%%%%%%%%%%%%%%%%%%%%%%%%%%%%%

\title{Agentes para Catan con heurísticas y modelos de lenguaje}
\author{Gómez Corral, Miquel \and Pachecho Millan, Andrés}
% \tutor{No}
% \curs{MUIARFID}

%%%%%%%%%%%%%%%%%%%%%%%%%%%%%%%%%%%%%%%%%%%%%%%%%%%%%%%%%%%%%%%%%%%%%%%%%%%%%%%
%                     PARAULES CLAU/PALABRAS CLAVE/KEY WORDS                  %
%                                                                             %
% Independentment de la llengua del treball, s'hi han d'incloure              %
% les paraules clau i el resum en els tres idiomes                            %
%%%%%%%%%%%%%%%%%%%%%%%%%%%%%%%%%%%%%%%%%%%%%%%%%%%%%%%%%%%%%%%%%%%%%%%%%%%%%%%

\keywords{
  Catan; agentes heurísticos; modelos de lenguaje; simulación multiagente; toma de decisiones
} % Paraules clau 
{
   Catan; agentes heurísticos; modelos de lenguaje; simulación multiagente; toma de decisiones
} % Palabras clave
{
  Catan; heuristic agents; language models; multi-agent simulation; decision making
} % Key words


%%%%%%%%%%%%%%%%%%%%%%%%%%%%%%%%%%%%%%%%%%%%%%%%%%%%%%%%%%%%%%%%%%%%%%%%%%%%%%%
%                                                                             %
%                              INICI DEL DOCUMENT                             %
%                                                                             %
%%%%%%%%%%%%%%%%%%%%%%%%%%%%%%%%%%%%%%%%%%%%%%%%%%%%%%%%%%%%%%%%%%%%%%%%%%%%%%%

\begin{document} 


%%%%%%%%%%%%%%%%%%%%%%%%%%%%%%%%%%%%%%%%%%%%%%%%%%%%%%%%%%%%%%%%%%%%%%%%%%%%%%%
%              RESUMENES DEL TFG EN VALENCIA, CASTELLA I ANGLES               %
%%%%%%%%%%%%%%%%%%%%%%%%%%%%%%%%%%%%%%%%%%%%%%%%%%%%%%%%%%%%%%%%%%%%%%%%%%%%%%%


% \begin{abstract}[spanish]

% \end{abstract}

%%%%%%%%%%%%%%%%%%%%%%%%%%%%%%%%%%%%%%%%%%%%%%%%%%%%%%%%%%%%%%%%%%%%%%%%%%%%%%%
%                                                                             %
%                              CONTENIDO DEL TREBAJO                          %
%                                                                             %
%%%%%%%%%%%%%%%%%%%%%%%%%%%%%%%%%%%%%%%%%%%%%%%%%%%%%%%%%%%%%%%%%%%%%%%%%%%%%%%

%%%%%%%%%%%%%%%%%%%%%%%%%%%%%%%%%%%%%%%%%%%%%%%%%%%%%%%%%%%%%%%%%%%%%%%%%%%%%%%
%                                  ÍDNICE                                     %
%%%%%%%%%%%%%%%%%%%%%%%%%%%%%%%%%%%%%%%%%%%%%%%%%%%%%%%%%%%%%%%%%%%%%%%%%%%%%%%
% \clearpage
% \tableofcontents
% \listoffigures
% \listoftables


\mainmatter

\chapter{Introducción}
En este trabajo contamos cómo hemos creado y evaluado un agente para \textit{PyCatan}. Hemos seguido dos caminos: reglas heurísticas (más rápidas y predecibles) y decisiones apoyadas por modelos de lenguaje (LLM). La idea principal es comprobar qué funciona mejor en cada tipo de decisión del juego y hasta qué punto compensa usar LLM en una simulación con muchas partidas.

\section{Alcance y estado actual}
Se han completado los experimentos de la parte heurística y se ha dejado preparada toda la parte LLM para tres proveedores (Ollama, AWS Bedrock y API UPV). A día de hoy tenemos:
\begin{itemize}
  \item Resultados completos del agente heurístico frente a random y estándar.
  \item Resultados \textbf{preliminares} de un \textit{smoke test} del agente LLM con Ollama.
  \item Integración lista para Bedrock y UPV, pero faltan pruebas completas porque aún no tenemos credenciales.
\end{itemize}

\chapter{Parte 1: Agente heurístico}

\section{Diseño del agente}
El agente heurístico se basa en reglas simples y rápidas de ejecutar. La ventaja es que su comportamiento es fácil de entender y depurar.

\subsection{Colocación inicial}
En \texttt{on\_game\_start}, el agente revisa los nodos válidos y da prioridad a:
\begin{itemize}
  \item Buena probabilidad total de producción en terrenos cercanos.
  \item Variedad de recursos para no quedarse atascado.
  \item Casillas potentes (6 y 8).
\end{itemize}

\subsection{Priorización de recursos y construcción}
En la fase de construcción sigue este orden:
\begin{enumerate}
  \item Construir ciudad si es posible.
  \item Construir poblado en el nodo con mejor potencial de producción.
  \item Construir carretera para mantener expansión.
  \item Comprar carta de desarrollo si no hay mejor opción estructural.
\end{enumerate}

También se corrigió el descarte por ladrón para evitar casos donde la partida podía quedarse bloqueada.

\section{Resultados del benchmark heurístico}
Con base en \texttt{results\_heuristic.json}:

\begin{table}[H]
\centering
\small
\begin{tabular}{lrrrrr}
  oprule
  extbf{Rival} & \textbf{Partidas} & \textbf{Victorias} & \textbf{Tasa de victoria} & \textbf{Puntos medios} & \textbf{Turnos medios} \\
\midrule
Random & 100 & 44 & 44.00\% & 7.07 & 444.65 \\
Estándar & 100 & 19 & 19.00\% & 5.00 & 281.80 \\
\bottomrule
\end{tabular}
\caption{Rendimiento del agente heurístico frente a random y agentes estándar.}
\label{tab:heuristic_results}
\end{table}

\paragraph{Lectura rápida}
El agente mejora claramente frente a random y compite de forma razonable frente a los agentes estándar. Aun así, hay margen de mejora, sobre todo en decisiones donde conviene adaptarse mejor al contexto de la partida.

\chapter{Parte 2: Integración de modelos de lenguaje}

\section{Cómo se conecta el agente con los LLM}
Se implementó una única interfaz para tres proveedores:
\begin{itemize}
  \item Ollama local.
  \item AWS Bedrock.
  \item API UPV.
\end{itemize}

La entrada y salida se definieron en JSON para evitar respuestas ambiguas. El agente LLM extiende al heurístico, así que si hay timeout o respuesta inválida, vuelve automáticamente al comportamiento heurístico.

\subsection{Decisiones que toma el LLM}
Se delegaron dos decisiones, como pide la práctica:
\begin{enumerate}
  \item \texttt{on\_game\_start}: selección de nodo y carretera inicial.
  \item \texttt{on\_build\_phase}: selección de acción de construcción (city/town/road/card/none).
\end{enumerate}

\subsection{Estado resumido para ahorrar tokens}
En vez de enviar toda la partida, enviamos solo lo importante:
\begin{itemize}
  \item Nodos válidos y resumen de terrenos cercanos (tipo y probabilidad) para la apertura.
  \item Recursos en mano y acciones legales para construir.
  \item Restricciones explícitas sobre formato de salida.
\end{itemize}

Con esto se reducen tokens, tarda menos y baja el número de respuestas mal formadas.

\section{Prompts diseñados (3 variantes)}

\subsection{Prompt A: formato estricto}
\begin{lstlisting}
Eres un asistente de Catan. Devuelve solo JSON valido.
Apertura: {"node_id": int, "road_to": int}
Construccion: {"building": "city|town|road|card|none", ...}
Elige solo acciones legales dentro de las opciones dadas.
\end{lstlisting}

\subsection{Prompt B: formato + validación interna}
\begin{lstlisting}
Piensa si la accion es legal y devuelve solo JSON valido.
No escribas texto fuera del JSON.
Si no hay accion legal, devuelve {"building": "none"}.
\end{lstlisting}

\subsection{Prompt C: formato + objetivo táctico}
\begin{lstlisting}
Objetivo: mejorar produccion a medio plazo y capacidad de expansion.
Prioriza acciones legales que aumenten ingresos de recursos.
Devuelve solo JSON, sin explicaciones.
\end{lstlisting}

\section{Resultados LLM y comparación de modelos}
	extbf{Estado de resultados:} de momento solo hay resultados completos de un \textit{smoke test} con Ollama. Bedrock y UPV están listos a nivel técnico, pero falta ejecutar pruebas completas cuando tengamos credenciales.

\subsection{Resultado preliminar disponible (Ollama)}
Con base en \texttt{results\_llm\_smoke.json}:

\begin{table}[H]
\centering
\small
\begin{tabular}{lrrrrr}
	oprule
	extbf{Rival} & \textbf{Partidas} & \textbf{Victorias} & \textbf{Tasa de victoria} & \textbf{Puntos medios} & \textbf{Turnos medios} \\
\midrule
Random & 4 & 1 & 25.00\% & 7.50 & 523.00 \\
Estándar & 4 & 3 & 75.00\% & 7.00 & 337.50 \\
\bottomrule
\end{tabular}
\caption{Resultado preliminar del agente LLM con Ollama (muestra pequeña).}
\label{tab:llm_smoke_results}
\end{table}

\subsection{Comparativa de modelos/proveedores (pendiente de completar)}
\begin{table}[H]
\centering
\small
\begin{tabular}{lcccc}
	oprule
	extbf{Proveedor/Modelo} & \textbf{Latencia media} & \textbf{Errores de formato} & \textbf{Calidad de decision} & \textbf{Estado} \\
\midrule
Ollama / llama3.2:3b & \textasciitilde 90 s por partida (smoke) & 0 en smoke & Media (muestra pequena) & Parcial \\
Bedrock / [PENDIENTE] & [MOCK] 0.0s & [MOCK] 0 & [MOCK] N/A & Falta ejecutar \\
UPV API / [PENDIENTE] & [MOCK] 0.0s & [MOCK] 0 & [MOCK] N/A & Falta ejecutar \\
\bottomrule
\end{tabular}
\caption{Comparativa provisional de modelos. Filas con [MOCK]/[PENDIENTE] deben reemplazarse tras experimentación real.}
\label{tab:llm_model_comparison}
\end{table}

\paragraph{Nota importante}
Los valores [MOCK] no son resultados reales. Están solo para dejar la tabla preparada y rellenarla cuando se hagan las pruebas pendientes.

\chapter{Parte 2.4: Análisis y conclusiones}

\section{Discusión}
La evidencia disponible confirma que:
\begin{itemize}
  \item El enfoque heurístico es estable, rápido y permite lanzar muchas partidas sin problemas.
  \item El enfoque LLM es más flexible, pero cada llamada tarda más.
  \item Si se simulan muchas partidas, ese tiempo extra se nota mucho.
\end{itemize}

\section{Viabilidad de LLM en este contexto}
El uso de LLM tiene más sentido en decisiones importantes que se toman pocas veces (por ejemplo, la apertura). Para decisiones repetitivas, sigue siendo mejor usar heurísticas o un enfoque híbrido.

\section{Qué falta para cerrar la parte experimental}
\begin{enumerate}
  \item Ejecutar benchmarks comparables con Bedrock y UPV en las mismas condiciones.
  \item Medir latencia por llamada y por partida para cada proveedor/modelo.
  \item Medir errores de formato JSON y tasa de fallback al agente heurístico.
  \item Comparar explícitamente Prompt A/B/C con igualdad de presupuesto de partidas.
\end{enumerate}

\chapter{Parte 3.1: Uso de herramientas de IA en el desarrollo}

Durante el desarrollo usamos herramientas de IA generativa para apoyar tareas de programación y redacción. En concreto, nos ayudaron en:
\begin{itemize}
  \item Automatizar scripts de benchmark con progreso y exportación de resultados.
  \item Diseñar una capa de abstracción de proveedores LLM (Ollama/Bedrock/UPV).
  \item Redactar borradores de secciones y tablas en LaTeX.
\end{itemize}

En general, estas herramientas aceleraron el trabajo y evitaron errores mecánicos. Aun así, las decisiones finales de diseño, la interpretación de resultados y la versión final del texto las hizo el equipo.

\section{Aportación humana pendiente de redactar}
	extbf{Sección para completar manualmente por los autores:}
\begin{itemize}
  \item Reparto exacto de tareas entre miembros del equipo.
  \item Valoración personal crítica sobre limitaciones encontradas con cada herramienta de IA.
  \item Decisiones de diseño que se descartaron y justificación final.
\end{itemize}

%%%%%%%%%%%%%%%%%%%%%%%%%%%%%%%%%%%%%%%%%%%%%%%%%%%%%%%%%%%%%%%%%%%%%%%%%%%%%%%
%                                                                             %
%                                BIBLIOGRAFIA                                 %
%                                                                             %
%%%%%%%%%%%%%%%%%%%%%%%%%%%%%%%%%%%%%%%%%%%%%%%%%%%%%%%%%%%%%%%%%%%%%%%%%%%%%%%
\cleardoublepage
\printbibliography

%%%%%%%%%%%%%%%%%%%%%%%%%%%%%%%%%%%%%%%%%%%%%%%%%%%%%%%%%%%%%%%%%%%%%%%%%%%%%%%
%                                                                             %
%                                 APÉNDICESS                                  %
%                                                                             %
%%%%%%%%%%%%%%%%%%%%%%%%%%%%%%%%%%%%%%%%%%%%%%%%%%%%%%%%%%%%%%%%%%%%%%%%%%%%%%%

\APPENDIX


%%%%%%%%%%%%%%%%%%%%%%%%%%%%%%%%%%%%%%%%%%%%%%%%%%%%%%%%%%%%%%%%%%%%%%%%%%%%%%
%                       EJEMPLOS DE CADA TIPO DE FACTURA                     %
%%%%%%%%%%%%%%%%%%%%%%%%%%%%%%%%%%%%%%%%%%%%%%%%%%%%%%%%%%%%%%%%%%%%%%%%%%%%%%

% \chapter{Apéndice correspondencia notebooks}
% \label{appendix:notebooks}
% ...

%%%%%%%%%%%%%%%%%%%%%%%%%%%%%%%%%%%%%%%%%%%%%%%%%%%%%%%%%%%%%%%%%%%%%%%%%%%%%%%
%                              FIN DEL DOCUMENTO                              %
%%%%%%%%%%%%%%%%%%%%%%%%%%%%%%%%%%%%%%%%%%%%%%%%%%%%%%%%%%%%%%%%%%%%%%%%%%%%%%%
\end{document}

